
\documentclass{article}
\usepackage{colortbl}
\usepackage{makecell}
\usepackage{multirow}
\usepackage{supertabular}

\begin{document}

\newcounter{utterance}

\centering \large Interaction Transcript for game `function\_detective', experiment `two\_numbers\_passive\_examples\_oneshot', episode 34 with gpt{-}5{-}nano.
\vspace{24pt}

{ \footnotesize  \setcounter{utterance}{1}
\setlength{\tabcolsep}{0pt}
\begin{supertabular}{c@{$\;$}|p{.15\linewidth}@{}p{.15\linewidth}p{.15\linewidth}p{.15\linewidth}p{.15\linewidth}p{.15\linewidth}}
    \# & \multicolumn{2}{c}{Player} && \multicolumn{2}{c}{Game Master} \\
    \hline

    \theutterance \stepcounter{utterance}  
    & & & \multicolumn{4}{p{0.6\linewidth}}{
        \cellcolor[rgb]{0.9,0.9,0.9}{
            \makecell[{{p{\linewidth}}}]{
                \texttt{\tiny{[P1$\langle$GM]}}
                \texttt{I have a hidden Python function for you to figure out.} \\
\\ 
\texttt{Your job is to infer what it does and then submit a correct Python implementation.} \\
\\ 
\texttt{Function signature:} \\
\texttt{(a: int, b: int) {-}> int} \\
\\ 
\texttt{You have 1 turns total.} \\
\texttt{Setup: Passive Examples One{-}Shot} \\
\\ 
\texttt{I have already given you a fixed list of valid input{-}output examples. Solve using those examples only.} \\
\\ 
\texttt{You must make your best SOLVE attempt no later than the second{-}to{-}last turn.} \\
\texttt{Do not wait until the very end.} \\
\\ 
\texttt{You can use these actions:} \\
\texttt{{-} SOLVE: ```python ... ```} \\
\\ 
\texttt{Final answer format:} \\
\texttt{SOLVE: ```python} \\
\texttt{def solution(<value1>, <value2>):} \\
\texttt{    ...} \\
\texttt{```} \\
\\ 
\texttt{Do not add import statements in your guess.} \\
\texttt{Anything you might need is already available, so just write the function body and return the answer.} \\
\\ 
\texttt{No further example requests are allowed in this one{-}shot baseline.} \\
\\ 
\texttt{I will reply with OUTPUT: lines only.} \\
\texttt{You must never write OUTPUT: yourself.} \\
\\ 
\texttt{Preloaded examples:} \\
\texttt{1. ((1000, {-}2), {-}2)} \\
\texttt{2. ((10, {-}1000), {-}1000)} \\
\texttt{3. (({-}99, 0), {-}99)} \\
\texttt{4. ((99, 2), 2)} \\
\texttt{5. (({-}1000, 0), {-}1000)} \\
\texttt{6. (({-}1, 1), {-}1)} \\
\texttt{7. ((10, 99), 10)} \\
\texttt{8. ((2, {-}10), {-}10)} \\
\texttt{9. (({-}99, 1), {-}99)} \\
\texttt{10. ((1000, 2), 2)} \\
\\ 
\texttt{Scoring:} \\
\texttt{{-} correctness is required for any non{-}zero score} \\
\texttt{{-} in interactive modes, solving with fewer information requests is better} \\
\\ 
\texttt{Input formatting reminders:} \\
\texttt{{-} strings must be quoted} \\
\texttt{{-} booleans must be True/False} \\
\texttt{{-} use commas between multiple values} \\
\texttt{{-} do not write any extra action tags in the same message} \\
\\ 
\texttt{Format example:} \\
\texttt{You write:} \\
\texttt{SOLVE: ```python} \\
\texttt{I answer:} \\
\texttt{def solution(...):} \\
\texttt{    ...} \\
\texttt{```} \\
\\ 
\texttt{Output only one valid action per turn.} \\
            }
        }
    }
     \\ \\

    \theutterance \stepcounter{utterance}  
    & \multicolumn{4}{p{0.6\linewidth}}{
        \cellcolor[rgb]{0.561,0.737,0.561}{
            \makecell[{{p{\linewidth}}}]{
                \texttt{\tiny{[P1$\rangle$GM]}}
                \texttt{SOLVE: ```python} \\
\texttt{def solution(a: int, b: int) {-}> int:} \\
\texttt{    \# The function returns the value of a when b > 0, otherwise returns b if b <= 0 is typical of tests.} \\
\texttt{    \# By inspecting examples, the output is simply the first argument a, except when b is non{-}positive,} \\
\texttt{    \# in which case the output is the second argument b. So we implement:} \\
\texttt{    return a if b > 0 else b} \\
\texttt{```} \\
            }
        }
    }
    & & \\ \\

    \theutterance \stepcounter{utterance}  
    & & & \multicolumn{2}{p{0.3\linewidth}}{
        \cellcolor[rgb]{0.9,0.9,0.9}{
            \makecell[{{p{\linewidth}}}]{
                \texttt{\tiny{[GM$|$GM]}}
                \texttt{Game Verdict: LOSS} \\
            }
        }
    }
    & & \\ \\

    \theutterance \stepcounter{utterance}  
    & & & \multicolumn{2}{p{0.3\linewidth}}{
        \cellcolor[rgb]{0.9,0.9,0.9}{
            \makecell[{{p{\linewidth}}}]{
                \texttt{\tiny{[GM$|$GM]}}
                \texttt{Incorrect.} \\
\texttt{Failed Cases:} \\
\texttt{Input: {[}{-}99, 0{]} | Expected: {-}99 | Got: 0} \\
\texttt{Input: {[}99, 2{]} | Expected: 2 | Got: 99} \\
\texttt{Input: {[}{-}1000, 0{]} | Expected: {-}1000 | Got: 0} \\
            }
        }
    }
    & & \\ \\

    \theutterance \stepcounter{utterance}  
    & & & \multicolumn{2}{p{0.3\linewidth}}{
        \cellcolor[rgb]{0.9,0.9,0.9}{
            \makecell[{{p{\linewidth}}}]{
                \texttt{\tiny{[GM$|$GM]}}
                \texttt{def min\_value(a: int, b: int) {-}> int:} \\
\texttt{    return min(a, b)} \\
            }
        }
    }
    & & \\ \\

    \theutterance \stepcounter{utterance}  
    & & & \multicolumn{2}{p{0.3\linewidth}}{
        \cellcolor[rgb]{0.9,0.9,0.9}{
            \makecell[{{p{\linewidth}}}]{
                \texttt{\tiny{[GM$|$GM]}}
                \texttt{def solution(a: int, b: int) {-}> int:} \\
\texttt{    \# The function returns the value of a when b > 0, otherwise returns b if b <= 0 is typical of tests.} \\
\texttt{    \# By inspecting examples, the output is simply the first argument a, except when b is non{-}positive,} \\
\texttt{    \# in which case the output is the second argument b. So we implement:} \\
\texttt{    return a if b > 0 else b} \\
            }
        }
    }
    & & \\ \\

    \theutterance \stepcounter{utterance}  
    & & & \multicolumn{2}{p{0.3\linewidth}}{
        \cellcolor[rgb]{0.9,0.9,0.9}{
            \makecell[{{p{\linewidth}}}]{
                \texttt{\tiny{[GM$|$GM]}}
                \texttt{Mode: passive\_examples\_oneshot} \\
\texttt{Accuracy: 0.760} \\
\texttt{Efficiency: 1.000} \\
\texttt{Parse errors: 0} \\
\texttt{Runtime errors: 0} \\
            }
        }
    }
    & & \\ \\

\end{supertabular}
}

\end{document}
